\section{Periodic weights and several distinct arithmetic poles}
\label{asj:sec:arithmetic}

\subsection{A complete quasi-polynomial specialization}
Let $q\ge1$ and let $a_0,\ldots,a_K$ be periodic functions on the positive
integers with period $q$. Define
\[
 c_n=\sum_{d=0}^K n^d a_d(n),\qquad
 \mu_d=\frac1q\sum_{j=1}^q a_d(j).
\]
Residue-class splitting in the absolute-convergence region gives
\begin{equation}\label{asj:eq:quasipoly}
 \D(s)=\sum_{d=0}^K q^{d-s}
                  \sum_{j=1}^q a_d(j)\zeta(s-d,j/q).
\end{equation}
This finite formula proves the continuation hypotheses of
Theorem~\ref{asj:thm:transfer}. Its only possible poles are simple, at $p=d+1$,
with residues $\mu_d$.

\begin{corollary}[Only the mean polynomial contributes local defects]
\label{asj:thm:mean}
For quasi-polynomial multiplicities and any polynomial $P$, all normalized
order-$r$ jets have degree at most $r+1$. For $r+1$ tangent vectors,
\begin{align}\label{asj:eq:mean}
 &\Delta_{\boldsymbol v_1}\cdots\Delta_{\boldsymbol v_{r+1}}
       \calJ_r(\boldsymbol\alpha)\notag\\
 &\quad=(-1)^{r+1}r![x^{-1}]P(x)
           \left(\sum_{d=0}^K\mu_dx^d\right)
                       \prod_i V_{\boldsymbol v_i}(x).
\end{align}
If every $\mu_d=0$, the degree is at most $r$ and the displayed difference
vanishes. More generally, replacing $c_n$ by its mean polynomial changes
only a degree-at-most-$r$ part of the normalized jet.
\end{corollary}
\begin{proof}
Use the residues in \eqref{asj:eq:quasipoly} in \eqref{asj:eq:top} and observe
$[x^{-(d+1)}]f(x)=[x^{-1}]x^d f(x)$. Subtracting the mean polynomial
leaves a base Dirichlet series whose possible poles have all been canceled,
so its shifted germ is holomorphic at the exponent origin. Theorem
\ref{asj:thm:polynomial} then gives degree at most $r$.
\end{proof}

For example, a nontrivial additive character $a(n)=e^{2\pi i kn/q}$ has
zero mean. Every normalized order-$r$ jet for this periodically weighted
spectrum, even after multiplication by $n^dP(n)$, has degree at most $r$.
There is no first-derivative multiplicative anomaly for such a pure
zero-mean periodic component. This does not say that its individual jets
vanish.

The finite parts themselves retain the periodic data. For one degree $d$,
the constant coefficient at $s=d+1$ of the corresponding term in
\eqref{asj:eq:quasipoly} is
\[
 \frac1q\sum_{j=1}^q a_d(j)\gamma_0(j/q)-\mu_d\log q.
\]
Other degrees contribute holomorphic values there. Higher coefficients
follow by multiplying $e^{-s\log q}$ by \eqref{asj:eq:stieltjes}. Thus the
mean law is a statement about the local defect, not the entire answer.

\subsection{A two-pole example from the sum-of-divisors function}
Let $\sigma_1(n)=\sum_{d\mid n}d$. The elementary product of two
absolutely convergent Dirichlet series proves
\begin{equation}\label{asj:eq:sigma1}
 \sum_{n\ge1}\frac{\sigma_1(n)}{n^s}=\zeta(s)\zeta(s-1),
 \qquad \Rea s>2.
\end{equation}
The two poles are simple, with
\[
 \kappa_{1,1}=\zeta(0)=-\frac12,
 \qquad \kappa_{2,1}=\zeta(2).
\]
The coexistence of distinct poles is genuinely different from inserting
a power of $\log n$ into one polynomially weighted spectrum.

For $P=1$, denote the continued multi-shift series by $Z_{\boldsymbol w}$,
and the single-shift series by $F_a$. Both are holomorphic at the spectral
origin, since all relevant numerators vanish there and the poles are simple.
Set
\[
 \delta_{\boldsymbol w}=
        \sum_j w_jF_{a_j}'(0)-Z_{\boldsymbol w}'(0),\qquad W\ne0.
\]
\begin{proposition}[An exact arithmetic anomaly]\label{asj:thm:sigma}
For $a_j>-1$,
\begin{equation}\label{asj:eq:sigma-anomaly}
 \boxed{\displaystyle
 \delta_{\boldsymbol w}
      =\frac{\zeta(2)}{2W}
                 \sum_{i<j}w_iw_j(a_i-a_j)^2.}
\end{equation}
\end{proposition}
\begin{proof}
Normalize the weights. For $r=1$, the nonlocal part is linear and cancels
against interpolation at the coordinate vertices. The pole at $1$ gives
$[x^{-1}]L_{\boldsymbol\alpha}^2=0$. The pole at $2$ gives
$\zeta(2)(\sum_j\alpha_ja_j)^2/2$. Therefore the difference between
vertex interpolation and the jet is
\[
 \frac{\zeta(2)}2\left\{
       \sum_j\alpha_ja_j^2-\left(\sum_j\alpha_ja_j\right)^2\right\}
   =\frac{\zeta(2)}2\sum_{i<j}\alpha_i\alpha_j(a_i-a_j)^2.
\]
Multiply by $W$ using \eqref{asj:eq:scaling}. The small-shift argument extends
to all stated shifts by the finite-head theorem.
\end{proof}

For two unit weights and shifts $a,b$, the anomaly is
$\zeta(2)(a-b)^2/4$. This is a concrete arithmetic specialization of the
classical simple-pole mechanism, not a claim to have discovered that
mechanism. In contrast, the next sections concern base functions with
higher poles, for which $F_a'(0)$ is typically not defined at all.
